\begin{lstlisting}[language=C, style = vscodecpp, caption=Newton-Raphson solver script for $A$ and $Z$, label=lst:nr]
#include <stdio.h>
#include <math.h>

// SEMF coefficients from [3]
double a1 = 15.56; // MeV
double a2 = 17.23;
double a3 =  0.697;
double a4 = 93.14;
double a5 = 12.00;

// Function definition f(A) - eq. 3.23
double f(double A) {
	double result =
		  2 * a2 * pow(a3,2) * pow(A,4.0/3.0)
		- a3 * pow(a4,2) * A 
		+ 4 * a2 * a3 * a4 * pow(A,2.0/3.0)
		+ 2 * a2 * pow(a4, 2);
	return result;
}

// Function derivative definition - fprime
double fp(double A) {
	double result = 
		  (8.0/3) * a2 * pow(a3,2) * pow(A, 1.0/3)
		- a3 * pow(a4,2)
		+ (8.0/3) * a2 * a3 * a4 * pow(A, -1.0/3);
	return result;
}

// Most stable Z for A
double Z(double A) {
	double num = a4 * A;
	double den = 2 * (a3 * pow(A,2.0/3) + a4);
	return num / den;
}

// Newton Raphson algorithm
double newton_raphson (double A0, double tol, int max_iter) 
{
	double A = A0;
	printf("Iteration 0: A = %.8f, f(A) = %12.5e\n", A, f(A));

	double f_val;  // function value
	double fp_val; // function derivative value
	double A_new;  // iterable var
	for (int i = 1; i <= max_iter; i++) {
		f_val = f(A);
		fp_val = fp(A);

		if (fabs(fp_val) < 1e-15) {
			printf("Derivative too small, method failure\n");
		}

		A_new = A - f_val / fp_val;
		printf("Iteration %d: A = %12.8f, f(A) = %12.5e\n", i, A_new, f(A_new));

		if (fabs(A_new - A) < tol) {
			printf("\nConverged after %d iterations.\n", i);
			return A_new;
		} 

		A = A_new;
	}

	printf("\nMax iterations reeached. Final A = %.8f\n", A);
	return A;
}

void horzline(char ch,int len) {
	for (int i = 0; i < len; i++) {
		printf("%c",ch);
	}
	printf("\n");
}
// ===============================
// MAIN ROUTINE
// ===============================
int main() 
{
	horzline('=',60);
	printf("Solving for A using NR method\n");
	horzline('=',60);

	double A0 = 100;
	printf("\nStarting with A0 = %.8f\n", A0);
	horzline('-',60);
	double solution = newton_raphson(A0, 1e-12, 100);
	printf("Maximal value for A: A = %.8f\n", solution);
	printf("Most stable element: Z = %.8f\n", Z(solution));
	horzline('=',60);
	return 0;
}
\end{lstlisting}
{\noindent \large \textbf{Output:}}
\begin{verbatim}
============================================================
Solving for A using NR method
============================================================

Starting with A0 = 100.00000000
------------------------------------------------------------
Iteration 0: A = 100.00000000, f(A) = -2.01545e+05
Iteration 1: A =  61.97467171, f(A) = -1.61385e+03
Iteration 2: A =  61.66458097, f(A) = -1.72475e-01
Iteration 3: A =  61.66454782, f(A) = -1.97906e-09
Iteration 4: A =  61.66454782, f(A) =  0.00000e+00

Converged after 4 iterations.
Maximal value for A: A = 61.66454782
Most stable element: Z = 27.40762339
============================================================
\end{verbatim}
